\section{System Under Test}
\label{sec:system}

The platform is load-bearing for this paper. A hybrid-graphics laptop with a
high-DPI panel under a scaling Wayland compositor supplies, by itself, the root
causes of a divide-by-zero in the display path (Section~\ref{sec:graphics1}), a
three-session misdiagnosis of the discrete GPU (Section~\ref{sec:egl}), a
window that appears absent from screenshots while existing (Section~\ref{sec:harness}),
and a game that renders in a box covering 40\% of the panel
(Section~\ref{sec:fix}). None of those would occur on a desktop with one GPU.

\subsection{Hardware}

\begin{table}[htbp]
\caption{Hardware under test}
\label{tab:hw}
\centering
\footnotesize
\begin{tabularx}{\columnwidth}{@{}lX@{}}
\toprule
\textbf{Component} & \textbf{Detail} \\
\midrule
Chassis & ASUS ROG Zephyrus G14, model GA403UU \\
Integrated GPU & AMD Radeon 780M (RDNA~3), no dedicated VRAM; borrows from
                 system memory via the GTT \\
Discrete GPU & NVIDIA GeForce RTX 4050 Laptop, 6141\,MiB dedicated VRAM,
               PCIe-attached, \emph{no attached display} \\
Graphics topology & NVIDIA Optimus (hybrid); the discrete GPU renders and
                    copies, it never scans out \\
System memory & 15.6\,GiB usable, shared bandwidth 102.4\,GB/s with the iGPU \\
Panel & 2880$\times$1800 physical \\
Storage & NVMe; game target on an NTFS volume with 699\,GB free \\
\bottomrule
\end{tabularx}
\end{table}

Two hardware facts recur throughout and are worth isolating now.

\textbf{The discrete GPU has no monitor attached.} On an Optimus system the panel
is wired to the integrated GPU. The discrete GPU is a compute-and-render device
whose output is blitted across PCIe. Any code path that enumerates display modes
for the rendering adapter therefore encounters an empty set, and code that
divides by a refresh rate obtained from that set faults
(Section~\ref{sec:graphics1}).

\textbf{The discrete GPU's memory is dedicated and non-elastic.} Unlike the
Radeon 780M, which borrows from system RAM and can therefore be over-committed
gracefully, the RTX 4050's 6141\,MiB is a hard budget. Anything already resident
on the card is subtracted from what the game can have. During this work a local
language-model service held 3576\,MiB of it, which is 58\% of the card. This did
not cause the central defect---we tested that explicitly and report the negative
result in Table~\ref{tab:eliminated-early}---but it shapes the final launcher
design in Section~\ref{sec:fix}.

\subsection{Software}

\begin{table}[htbp]
\caption{Software stack}
\label{tab:sw}
\centering
\footnotesize
\begin{tabularx}{\columnwidth}{@{}lX@{}}
\toprule
\textbf{Layer} & \textbf{Version / detail} \\
\midrule
OS & Arch Linux, rolling \\
Display & KDE Plasma on Wayland, KWin; scale factor 2 \\
X11 clients & XWayland, \emph{unscaled} \\
NVIDIA driver & 595.71.05, upgraded to 610.57.04 during the work \\
Mesa & 26.1.6 (RADV for the Radeon 780M) \\
Compatibility layer & GE-Proton~10-34; system Wine 11.14 for contrast \\
D3D12 translation & vkd3d-proton (bundled), 2.14.1, and 3.0.1 \\
D3D11/10/9 translation & DXVK \\
Vulkan loader & ICD selection via \verb|VK_ICD_FILENAMES| \\
EGL dispatch & libglvnd, vendor list pinned system-wide \\
\bottomrule
\end{tabularx}
\end{table}

\subsection{The scaling discontinuity}

KDE runs at scale factor 2, so the 2880$\times$1800 panel presents as
1440$\times$900 logical pixels to Wayland-native clients. \textbf{XWayland is
unscaled}, so X11 clients---which includes everything under Wine---work in
physical pixels and see a 2880$\times$1800 display.

This produces two distinct and independently confusing artefacts:

\begin{enumerate}
\item \textbf{Screenshots and compositor coordinates are different number
      systems.} A screenshot is 2880$\times$1800; the compositor's own cursor and
      geometry queries answer in 1440$\times$900. Every pixel measured off a
      screenshot must be halved before it can be used to aim anything. Nothing
      warns about this; the symptom is aiming consistently twice as far right and
      down as intended.
\item \textbf{A correct window looks half-size in compositor reports.} During the
      deadlock investigation, KWin described the game window as
      \verb|960x568|, which reads as a scaling bug. It is correct: a
      1920$\times$1080 physical window is 960$\times$540 logical, plus a 28-pixel
      titlebar. The window was the right size the entire time.
\end{enumerate}

\subsection{The discrete-GPU access gate}

The system enforces a default-deny policy on the discrete GPU. The device node
\verb|/dev/nvidia0| is owned \verb|root:dgpu| with mode 0660, and the
\verb|dgpu| group is deliberately empty. Access is granted only through a setgid
helper, \verb|dgpu-exec-v2|, mode \verb|-rwxr-sr-x root dgpu|. The policy exists
so that the discrete GPU can reach its runtime power-management sleep state
instead of being woken by incidental enumeration.

Two consequences matter for the investigation.

First, \textbf{the gate is the intuitive suspect for every discrete-GPU failure
and was never once the cause}. We record this because it was the first
hypothesis on three separate occasions. It is falsifiable cheaply and in a way
worth reusing: the failures reproduce identically \emph{as root} and with the
gate open, which removes permissions from the space of explanations. In the
central defect the symptom shape also rules it out on its own---a gate failure
produces zero dialogs and a silent fallback to the integrated GPU, whereas
reaching a Vulkan pipeline-creation call at all proves the device was already
opened successfully.

Second, \textbf{a setgid process is non-dumpable}. Anything launched through the
helper has \verb|/proc/PID/mem|, \verb|/proc/PID/maps|, and
\verb|/proc/PID/environ| readable only by root, even for the invoking user.
\verb|/proc/PID/status| remains readable, which is what made the gate's operation
verifiable at all: the game's status file reported
\verb|Gid: 948 948 948 948|, confirming the setgid transition had taken effect.

A related defect found during this work is worth one line, because it is a
general hazard in privilege-raising wrappers: an earlier version of the helper
raised only the \emph{effective} group ID, and \verb|bash| drops that on startup.
Any wrapper invoked through a shell script therefore silently lost the privilege
it had just acquired. The launcher must invoke the v2 helper, which raises the
real GID as well, and it must exec the target binary rather than a shell.

\subsection{Nomenclature used throughout}

\begin{table}[htbp]
\caption{Terms and abbreviations used in this paper}
\label{tab:terms}
\centering
\footnotesize
\begin{tabularx}{\columnwidth}{@{}lX@{}}
\toprule
\textbf{Term} & \textbf{Meaning here} \\
\midrule
\emph{box} & one of the six \verb|.bin| containers \\
\emph{prefix} & a Wine or Proton emulated Windows filesystem and registry \\
\emph{ICD} & Installable Client Driver; a vendor's Vulkan implementation \\
DXIL & the shader bytecode format Direct3D~12 consumes \\
SPIR-V & the shader bytecode format Vulkan consumes \\
\assertn{n} & count of Win32 dialog windows found in the prefix; the
              failure metric defined in Section~\ref{sec:harness} \\
\verb|#32770| & the Win32 window class of a dialog box \\
\bottomrule
\end{tabularx}
\end{table}
